\section{Building and Releasing}
\label{sec:build}

The repository holds three kinds of files: the class, the content, and the tooling.
\Cref{fig:arch} already showed how they connect; what matters here is that moving to a
new machine should cost no thinking.

\subsection{Local builds}

On Windows use \texttt{build.ps1}; on macOS and Linux use the \texttt{Makefile}. Both
drive \texttt{latexmk} through two passes and the bibliography. The script picks the
engine from the language, so you never have to remember which file needs XeLaTeX.

\begin{lstlisting}[language=bash,
  caption={Common build commands. With \texttt{-Job}, only one report is built.},
  label={lst:build}]
./build.ps1                 # build both example reports
./build.ps1 -Job main-en    # English only
./build.ps1 -Clean          # remove intermediates
./build.ps1 -Watch          # rebuild on save
make en                     # macOS / Linux
\end{lstlisting}

\subsection{Continuous integration}

CI does two things: it proves that a clean machine can build the PDF, and it keeps a copy
of the result. Two jobs run in parallel, English on pdfLaTeX and Chinese on XeLaTeX with
the Fandol character set. Every push builds and uploads the PDF as an artifact; a tag also
creates a release with both reports attached.

\begin{lstlisting}[language=yaml,
  caption={The essential CI fragment. The Chinese job injects \texttt{fandol} through \texttt{\textbackslash PassOptionsToClass} so it never depends on host fonts.},
  label={lst:ci}]
jobs:
  build:
    strategy:
      matrix:
        include:
          - job: main-en
          - job: main-zh
    steps:
      - uses: actions/checkout@v4
      - uses: xu-cheng/latex-action@v3
        with:
          root_file: ${{ matrix.job }}.tex
\end{lstlisting}

CI proves compilation, nothing more. It cannot see an overfull table or a figure stranded
on its own page; layout still needs a human pass, which is why the release checklist
exists.

\subsection{Versions}

The version lives in \verb|\reportsetup{version={...}}| and matches the git tag. Bump it
together with the visual check, the changelog and the release notes. Done separately,
any one of those three becomes a mystery six months later.
